% !PRINTED 164
\weberpage{164}
Hat die Gruppe $S$ vom Grade $\mu$ eine Invariante vom Index
 $\mu$, so wird $T$ die Einheitsgruppe und $S$ ist eine cyklische
Gruppe, die aus den Elementen $1, E, E^2 \ldots E^{\mu-1}$ besteht.

Betrachten wir als Beispiel die symmetrische Permutationsgruppe
 $P$ von $n$ Elementen $x_1,x_2,\ldots,x_n$, die ja nach \S.\,39 unter
den Gruppen linearer Substitutionen enthalten ist, so haben wir
als absolute Invarianten die \gesperrt{symmetrischen Functionen} der
$n$ Variablen $x$. Das Differenzenproduct
\[
\sqrt{\varDelta} = (x_1-x_2)(x_1-x_3)\ldots(x_{n-1}-x_n)
\]
ist eine relative Invariante vom Index 2, zu der die alternirende
Gruppe gehört. Da die Gruppe $P$ ausser der alternirenden
Gruppe keinen Normaltheiler hat und auch nicht cyklisch ist, so
giebt es keine relativen Invarianten von höherem Index als 2.

Die absoluten Invarianten, d. h. die symmetrischen Functionen
, sind hier durch eine endliche Anzahl solcher Formen
rational darstellbar, nämlich durch die symmetrischen Grundfunctionen
, und die Invarianten vom Index 2 sind das Product
von $\sqrt{\varDelta}$ mit absoluten Invarianten. Dass analoge S\"atze auch
im allgemeinen Falle gelten, werden wir in der Folge beweisen.

Wir schliessen hier mit dem Beweise eines allgemeinen
Satzes, der bei allen Anwendungen f\"ur die Bildung der Invarianten
einer Gruppe $S$ von grossem Nutzen ist.

Im \S.\,60 des ersten Bandes haben wir f\"ur irgend eine Form
der $n$ Variablen $F(x_1,x_2,\ldots x_n) = F(x)$ gewisse Formen derselben
 Variablen $C(x_1,x_2,\ldots x_n) = C(x)$ als Covarianten
definirt, die dadurch charakterisirt waren, dass, wenn $F(x)$ durch
irgend eine lineare Substitution in eine neue Form $F'(y)$ transformirt
 wird, die Form $C$ der Bedingung genügt
\[
C'(y) = r^\lambda\,C(x),
\]
wo $C'$ ebenso von den Co\"efficienten von $F'$, wie $C$ von den
Co\"efficienten von $F$ abhängt. Darin bedeutet $r$ die Substitutionsdeterminante
, ist also eine Constante. Die Co\"efficienten der
Form $F$ kommen in $C$ nur homogen vor. Als Beispiel einer Covariante
 f\"uhren wir die \gesperrt{Hesse'sche Determinante} an.

Wenn nun $F(x)$ eine Invariante der Gruppe $S$ ist, und die
$y$ mit den $x$ gleichfalls durch eine Substitution aus $S$ zusammen-
h\"angen, so unterscheiden sich die Co\"efficienten von $F$ nur durch
einen gemeinschaftlichen constanten Factor von den Co\"efficienten
von $F'$, und Gleiches gilt also auch von den beiden Formen $C$
%CONT>>
